\chapter{Vaginitis}

\section*{Definition}
Vaginitis is inflammation of the vaginal mucosa presenting with abnormal discharge, odor, itching, burning, or irritation. The three most common etiologies are bacterial vaginosis, vulvovaginal candidiasis, and trichomoniasis, each with distinct clinical and microbiological features.

\section*{What Happens in Your Body}
Vaginitis results from disruption of the normal protective vaginal microbiome (dominant Lactobacillus species that produce hydrogen peroxide and maintain an acidic pH of 3.8--4.5). Pathogen overgrowth (bacteria, yeast, or protozoa) triggers an inflammatory response via epithelial immune signaling proteins release and recruitment of neutrophils.

\section*{Causes}
\begin{itemize}
  \item \textbf{Bacterial vaginosis (BV):} Overgrowth of Gardnerella vaginalis, Atopobium vaginae, and other anaerobes; thin, gray/white discharge with fishy odor; elevated pH >4.5
  \item \textbf{Vulvovaginal candidiasis:} Candida albicans (80--90\%) or non-albicans species; thick, white, cottage-cheese discharge; itching and burning; normal pH 4--4.5
  \item \textbf{Trichomoniasis:} Trichomonas vaginalis (sexually transmitted); green-yellow, frothy discharge; itching and painful urination; elevated pH >5; motile organisms on wet mount
  \item \textbf{Atrophic vaginitis:} Estrogen deficiency (postmenopausal, breastfeeding, postpartum); thin, inflamed mucosa; pain during intercourse; dry or watery discharge
  \item \textbf{Irritant/allergic:} Spermicides, douches, soaps, fragrances, latex, panty liners
  \item \textbf{Other infections:} Group B Streptococcus, Mycoplasma, Chlamydia, gonorrhea, HSV
\end{itemize}

\section*{When to See a Doctor}
See your doctor if:
\begin{itemize}
  \item Discharge is new, abnormal, or associated with itching, burning, or odor
  \item Suspected sexually transmitted infection (trichomoniasis, cervicitis)
  \item Pregnant or postoperative (vaginitis carries pregnancy and surgical risks)
  \item Recurrent episodes (4 or more per year)
\end{itemize}

\section*{Related Symptoms}
\begin{itemize}
  \item Abnormal vaginal discharge, itching, vulvar irritation, painful urination, pain during intercourse, odor, vaginal redness, vaginal dryness
\end{itemize}
\textbf{Important:} Bacterial vaginosis is the most common cause of vaginitis in reproductive-age women, affecting 15--50\% depending on the population.

\section*{Self-Care / Home Management}
\begin{itemize}
  \item Avoid douching, scented hygiene products, and irritants
  \item Wear loose, cotton underwear and change out of wet clothing promptly
  \item Apply cold compresses for vulvar itching and irritation
  \item Over-the-counter antifungal (clotrimazole, miconazole) for uncomplicated candidiasis
\end{itemize}

\section*{How Your Doctor Will Evaluate You}
\begin{itemize}
  \item Vaginal pH testing (normal 3.8--4.5): pH >4.5 suggests BV or trichomoniasis
  \item Wet mount microscopy: clue cells (BV), pseudohyphae (Candida), motile trichomonads (Trichomonas)
  \item Whiff test (amine odor with KOH): positive in BV
  \item Nucleic acid amplification test (NAAT) for Trichomonas and other STIs
  \item Culture if resistant or recurrent Candida is suspected
\end{itemize}

\section*{Treatment Options}
\begin{itemize}
  \item \textbf{BV:} Metronidazole (oral or topical) or clindamycin cream; treat sexual partner not required
  \item \textbf{Candidiasis:} Topical azoles (clotrimazole, miconazole, terconazole) or oral fluconazole
  \item \textbf{Trichomoniasis:} Oral metronidazole or tinidazole; treat all sexual partners
  \item \textbf{Atrophic:} Topical vaginal estrogen cream, ring, or tablet
  \item Recurrent BV or candidiasis: maintenance suppressive therapy and evaluation for underlying risk factors (diabetes, immunosuppression)
\end{itemize}

\section*{References}
\begin{thebibliography}{99}
\bibitem{vaginitis:ref1} Sobel JD. ``Vulvovaginal candidosis.'' \textit{The Lancet}, 2007;369(9577):1961--1971.
\bibitem{vaginitis:ref2} Workowski KA, Bachmann LH, Chan PA, et al. ``Sexually Transmitted Infections Treatment Guidelines, 2021.'' \textit{MMWR Recommendations and Reports}, 2021;70(4):1--187.
\bibitem{vaginitis:ref3} Nyirjesy P. ``Management of persistent vaginitis.'' \textit{Obstetrics \& Gynecology}, 2014;124(6):1135--1146.
\bibitem{vaginitis:ref4} Hillier SL, Holmes KK. ``Bacterial vaginosis.'' In: \textit{Sexually Transmitted Diseases}, 4th ed. McGraw-Hill, 2008.
\bibitem{vaginitis:ref5} Schwebke JR, Muzny CA, Josey WE. ``Role of Gardnerella vaginalis in the pathogenesis of bacterial vaginosis: A conceptual model.'' \textit{Journal of Infectious Diseases}, 2014;210(3):338--343.
\bibitem{vaginitis:ref6} Hainer BL, Gibson MV. ``Vaginitis: Diagnosis and treatment.'' \textit{American Family Physician}, 2011;83(7):807--815.
\end{thebibliography}
